\section{From affine labels to persistent trajectories}
\label{sec:dynamics}

We now apply Lemma~\ref{lem:affine} to the mass-action vector field.
Its uniform form in Corollary~\ref{cor:uniform-affine} lets us choose
the tolerance so that every strict ordering made by an active slope gives
a monomial comparison with a definite gap. Weak reversibility converts those comparisons into
nonnegative derivatives, and an initially active constant branch supplies
the compact set that traps the trajectory.

\subsection{Uniform approximation and sign constraints}

Fix henceforth a network, positive rate constants, and
$x^0\in\R_{>0}^d$ as in Theorem~\ref{thm:main}. Write $e_1,\ldots,e_d$
for the standard basis of $\R^d$, and set
\begin{equation}\label{eq:comparison-directions}
 D=\{e_1,\ldots,e_d\}
 \cup\{z-y:y,z\in\mathcal C,\ y\ne z\}.
\end{equation}
Every vector in $D$ is nonzero. We include differences of \emph{all}
pairs of complexes, not only reaction vectors. The flux comparison below
compares the source of one upward reaction with the sources of every
downward reaction. Such sources need not be adjacent in the reaction
graph and may even belong to different linkage classes.

For $r\in[-1,1]^d$, define
\begin{equation}\label{eq:chosen-tolerance}
 E(r)=\min\left(
 \left\{\frac12\right\}
 \cup
 \left\{
 \frac{\abs{r\cdot v}}{2\norm{v}_1}:
 v\in D,\ r\cdot v\ne0
 \right\}\right),
 \qquad \norm{v}_1=\sum_{i=1}^d\abs{v_i}.
\end{equation}
This is a positive number for every $r$, since the displayed minimum is
over a finite nonempty set of positive numbers; in particular,
$E(0)=1/2$. The definition is allowed to be discontinuous at the
hyperplanes $r\cdot v=0$. This is the reason for allowing arbitrary
positive tolerances in Lemma~\ref{lem:affine}.

Apply that Lemma with $m=d$, $a=-1$, $b=1$, and the tolerance
\eqref{eq:chosen-tolerance}. Denote its fixed family by
$\{(r_\lambda,c_\lambda):\lambda\in\Lambda\}$ and put
\[
 F_h(x)=\min_{\lambda\in\Lambda}
       \bigl(r_\lambda\cdot x+c_\lambda(h)\bigr)
 \qquad(x\in\R^d).
\]
Choose $h_0\in(0,1)$ small enough that all offsets are defined and
Corollary~\ref{cor:uniform-affine} applies for $0<h<h_0$.
Before choosing a particular $h$, define
\begin{equation}\label{eq:fixed-gap}
 H=\min\left(
 \{1\}\cup
 \left\{
 \frac12\abs{r_\lambda\cdot v}:
 \lambda\in\Lambda,\ v\in D,\ r_\lambda\cdot v\ne0
 \right\}\right)>0.
\end{equation}
The positivity of $H$ again follows from finiteness. Both the slopes and
$H$ are independent of the eventual choice of $h$.
Thus the order of choices is
\[
 D\ \longrightarrow\ E\ \longrightarrow\ \Lambda
 \ \longrightarrow\ (h_0,H)\ \longrightarrow\ h.
\]
In particular, $H$ is not a separation parameter that is allowed to tend
to zero together with $h$.

Suppose that $0<h<h_0$, $p\in[-1,1]^d$, and $\lambda$ is active at
$X_h(p)$. For every $v\in D$ with $r_\lambda\cdot v\ne0$, we have
\begin{align*}
 \abs{(p-r_\lambda)\cdot v}
 &\le \norm{p-r_\lambda}_\infty\norm{v}_1\\
 &<E(r_\lambda)\norm{v}_1
 \le\frac12\abs{r_\lambda\cdot v}.
\end{align*}
Consequently,
\begin{equation}\label{eq:sign-and-gap}
 \operatorname{sgn}(p\cdot v)
 =\operatorname{sgn}(r_\lambda\cdot v),
 \qquad \abs{p\cdot v}\ge H.
\end{equation}
In particular, $\abs{p\cdot v}<H$ forces $r_\lambda\cdot v=0$.
Thus a comparison too close to a tie in exponent space is made an exact
tie by the active slope. The converse is not required: equal slope levels
may contain complexes with different monomial values.
For $x=X_h(p)$ and any $y,z\in\mathcal C$, this implies
\begin{equation}\label{eq:monomial-order}
 r_\lambda\cdot y<r_\lambda\cdot z
 \quad\Longrightarrow\quad
 x^z\le h^H x^y.
\end{equation}
Indeed, $z-y\in D$ and $p\cdot(z-y)\ge H$, while
$x^z/x^y=h^{p\cdot(z-y)}$ and $0<h<1$.

We can now fix a single $h\in(0,h_0)$ satisfying
\begin{equation}\label{eq:initial-exponents}
 p_i^0=\frac{\log x_i^0}{\log h},
 \qquad \norm{p^0}_\infty<H,
\end{equation}
and, when $\mathcal R\ne\varnothing$, also satisfying
\begin{equation}\label{eq:rate-separation}
 h^H\sum_{e\in\mathcal R}\kappa_e
 <\min_{e\in\mathcal R}\kappa_e.
\end{equation}
This choice is possible because $p^0\to0$ and $h^H\to0$ as
$h\to0+$, with $H>0$ already fixed. No minimum of reaction rates is
needed when $\mathcal R$ is empty. Write
\[
 B_h=[h,h^{-1}]^d,
 \qquad F=F_h.
\]
Since $H\le1$, Equation~\eqref{eq:initial-exponents} places $x^0$ in
the interior of $B_h$.

Every label active at $x^0$ has slope zero. In fact, if its $i$th slope
coordinate were nonzero, Equation~\eqref{eq:sign-and-gap} applied to
$v=e_i$ would give $\abs{p_i^0}\ge H$, contradicting
\eqref{eq:initial-exponents}. The finite nonempty family has at least
one active label; fix such a label $\lambda_0$ and put $M=F(x^0)$. Then
\begin{equation}\label{eq:ceiling}
 r_{\lambda_0}=0,
 \qquad M=F(x^0)=c_{\lambda_0}(h),
 \qquad F(x)\le M\quad(x\in\R^d).
\end{equation}
Thus a constant branch is forced by the sign constraints, even though it
was not prescribed in the application of Lemma~\ref{lem:affine}.

There is no point $x\in\partial B_h$ at which $F(x)=M$. To see this,
write $x=X_h(p)$ with $p\in[-1,1]^d$. The boundary condition is
$\norm{p}_\infty=1$. Equality $F(x)=M$ would make $\lambda_0$ active,
so Corollary~\ref{cor:uniform-affine} would instead give
$\norm{p}_\infty<E(0)=1/2$. We have proved
\begin{equation}\label{eq:no-boundary-plateau}
 \{x\in B_h:F(x)=M\}\cap\partial B_h=\varnothing.
\end{equation}
Thus the maximum set $\{F=M\}$ contains $x^0$ and misses
$\partial B_h$. To prove its invariance, we next show that a trajectory
segment in $B_h$ cannot decrease $F$. The constant branch alone does not prove this: another branch tied
with it might otherwise become smaller. We must control the derivatives
of all active branches. The next two arguments establish that control
from weak reversibility.

\subsection{Flux across the ordered complex levels}

The following finite graph statement isolates the role of weak
reversibility. Its weights will be the monomials $x^y$, and its levels
will be the scalar products $r_\lambda\cdot y$.

\begin{lemma}[Flux across level cuts]\label{lem:cut-flux}
Let $V$ be a finite set and $\mathcal E$ a finite set of directed edges
with source and target maps $s,t:\mathcal E\to V$. Suppose that for
every edge there is a directed path from its target back to its source.
Let $a:V\to\R$, let $w_v>0$ for $v\in V$, and let
$\kappa_e>0$ for $e\in\mathcal E$. Assume that there is $\delta>0$
such that
\begin{equation}\label{eq:abstract-weight-order}
 a(u)<a(z)\quad\Longrightarrow\quad w_z\le\delta w_u
 \qquad(u,z\in V).
\end{equation}
If $\mathcal E\ne\varnothing$, assume in addition that
\[
 \delta\sum_{e\in\mathcal E}\kappa_e
 <\min_{e\in\mathcal E}\kappa_e.
\]
Then
\[
 \sum_{e\in\mathcal E}
 \kappa_e w_{s(e)}\bigl(a(t(e))-a(s(e))\bigr)\ge0.
\]
\end{lemma}

\begin{proof}
The sum is zero if $\mathcal E$ is empty, including the case of an empty
vertex set. Otherwise list the distinct values of $a$ as
$\alpha_1<\cdots<\alpha_q$. For $1\le j<q$, let
\[
 V_j^- =\{v\in V:a(v)\le\alpha_j\},
 \qquad V_j^+=V\setminus V_j^-.
\]
Let $\mathcal E_j^+$ consist of edges from $V_j^-$ to $V_j^+$ and
$\mathcal E_j^-$ of edges in the reverse direction. Define their net
upward flux by
\[
 Q_j=\sum_{e\in\mathcal E_j^+}\kappa_e w_{s(e)}
      -\sum_{e\in\mathcal E_j^-}\kappa_e w_{s(e)}.
\]
If $\mathcal E_j^-\ne\varnothing$, take one of its edges. Its return
path starts in $V_j^-$ and ends in $V_j^+$, so that path contains an edge
$e_0\in\mathcal E_j^+$. Put $u=s(e_0)$. Every source $z$ of an edge
in $\mathcal E_j^-$ satisfies $a(z)>a(u)$, and hence
\begin{align*}
 \sum_{e\in\mathcal E_j^-}\kappa_e w_{s(e)}
 &\le \delta w_u\sum_{e\in\mathcal E_j^-}\kappa_e\\
 &\le \delta w_u\sum_{e\in\mathcal E}\kappa_e
 <w_u\min_{e\in\mathcal E}\kappa_e
 \le \kappa_{e_0}w_u.
\end{align*}
It follows that $Q_j>0$. If $\mathcal E_j^-$ is empty, then $Q_j\ge0$
directly.
Notice that the selected edge $e_0$ need not share a component with every
downward edge. The all-pairs hypothesis
\eqref{eq:abstract-weight-order} is what permits one upward contribution
to dominate the entire downward sum in this proof.

For each edge, its level difference is the signed sum of the gaps
$\alpha_{j+1}-\alpha_j$ over the cuts it crosses. Summing this identity
over the finite edge set gives
\begin{equation}\label{eq:cut-decomposition}
 \sum_{e\in\mathcal E}
 \kappa_e w_{s(e)}\bigl(a(t(e))-a(s(e))\bigr)
 =\sum_{j=1}^{q-1}(\alpha_{j+1}-\alpha_j)Q_j\ge0.
\end{equation}
When $q=1$, both sides vanish. Equal vertex levels therefore cause no
difficulty, and no connectedness assumption has entered the argument.
\end{proof}

Return-path monomial domination appears in Anderson's tier
argument~\cite[Lemma~4.7]{Anderson2011Global}, and quantitative inward
reaction estimates underlie the invariant polygons of Craciun, Nazarov,
and Pantea~\cite{CraciunNazarovPantea2013}. Ordered-cut regrouping also
appears in the proof of Craciun's toric-inclusion embedding
theorem~\cite[preprint, Theorem~4.1]{Craciun2019}. Lemma~\ref{lem:cut-flux}
uses a direct finite estimate: one upward edge dominates the full downward
flux. No tier extraction, cycle decomposition, or differential inclusion
is required in the proof above.

For any $x\in B_h$ and any label $\lambda$ active at $x$, apply
Lemma~\ref{lem:cut-flux} with
\[
 V=\mathcal C,\quad \mathcal E=\mathcal R,\quad
 a(y)=r_\lambda\cdot y,\quad w_y=x^y,\quad \delta=h^H.
\]
The return paths are supplied by weak reversibility,
Equation~\eqref{eq:monomial-order} supplies
\eqref{eq:abstract-weight-order}, and
Equation~\eqref{eq:rate-separation} supplies the rate inequality when
there are reactions. Thus
\begin{equation}\label{eq:active-nondecrease}
 r_\lambda\cdot f(x)
 =\sum_{y\to y'\in\mathcal R}
   \kappa_{y\to y'}x^y\,r_\lambda\cdot(y'-y)
 \ge0
 \qquad(x\in B_h,\ \lambda\text{ active at }x).
\end{equation}
This includes $f=0$ for an empty reaction set. The comparison in
\eqref{eq:abstract-weight-order} is required for all ordered vertex
pairs, including nonadjacent pairs in one component and pairs in different
components; this explains the full set of complex differences in
\eqref{eq:comparison-directions}.

\subsection{Derivatives of finite affine minima}

Although $F$ need not be differentiable where its active labels change,
its right derivative along a differentiable curve is explicit. We record
the finite-affine minimum case of the classical active-branch
directional-derivative principle; see, for example,
Bertsekas~\cite[Proposition~A.3.2(a)]{Bertsekas2009} for the corresponding
maximum formula. The short proof includes all ties and the monotonicity
implication needed here.

\begin{lemma}[Right derivative of a finite affine minimum]
\label{lem:minimum-derivative}
Let $I$ be a finite nonempty index set, let $b_i\in\R^d$ and
$d_i\in\R$ for $i\in I$, and define
\[
 \Phi(x)=\min_{i\in I}(b_i\cdot x+d_i),
 \qquad A(x)=\{i\in I:b_i\cdot x+d_i=\Phi(x)\}.
\]
For a continuously differentiable curve
$\gamma:[u,v]\to\R^d$, the right derivative at every $\tau\in[u,v)$
exists and is
\[
 \frac{d^+}{d\tau}\Phi(\gamma(\tau))
 =\min_{i\in A(\gamma(\tau))} b_i\cdot\gamma'(\tau).
\]
If $b_i\cdot\gamma'(\tau)\ge0$ for every $\tau\in[u,v)$ and every
$i\in A(\gamma(\tau))$, then $\Phi\circ\gamma$ is nondecreasing on
$[u,v]$.
\end{lemma}

\begin{proof}
Fix $\tau<v$. Each inactive branch has a strictly positive gap above
any chosen active branch at $\gamma(\tau)$. By continuity and finiteness,
every inactive branch stays above that chosen branch at
$\gamma(\tau+\eta)$ for sufficiently small $\eta\ge0$. For the remaining
branches, the first-order expansions, uniform over the finite active set,
give
\[
 \Phi(\gamma(\tau+\eta))
 =\Phi(\gamma(\tau))
  +\eta\min_{i\in A(\gamma(\tau))}
       b_i\cdot\gamma'(\tau)+o(\eta)
 \qquad(\eta\downarrow0).
\]
This proves the right derivative formula, including at ties.

Under the stated nonnegativity hypothesis, take any
$u\le s<t\le v$ and any $\epsilon>0$. The continuous function
\[
 g(\tau)=\Phi(\gamma(\tau))+\epsilon\tau
 \qquad(s\le\tau\le t)
\]
has strictly positive right derivative at every $\tau<t$. It cannot
attain its maximum at such a point, because sufficiently small positive
increments would increase its value. Compactness therefore places its
maximum at $t$, giving $g(s)\le g(t)$. Letting $\epsilon\downarrow0$
proves $\Phi(\gamma(s))\le\Phi(\gamma(t))$.
\end{proof}

\subsection{The invariant polytope and trajectory bounds}

All active branches have nonnegative derivative in $B_h$, while the
constant branch in \eqref{eq:ceiling} gives $F\le M$ everywhere. We now
show that the maximum set of $F$ is the compact invariant polytope
announced in the introduction. Its invariance will give
Theorem~\ref{thm:main}.

Nonnegative derivatives on all active inequalities express the familiar
inward-tangency condition for invariance, whose classical general form
goes back to Nagumo~\cite{Nagumo1942}. Here we prove invariance directly
by a first-exit argument.

\begin{proposition}\label{prop:polytope}
For every finite weakly reversible mass-action network with positive
constant rates, and every $x^0\in\R_{>0}^d$, there is a compact convex
polytope $K\subset\R_{>0}^d$ containing $x^0$ and forward invariant for
the mass-action equation. Here forward invariance means that every
solution starting in $K$ exists for all forward time and stays in $K$.
\end{proposition}

\begin{proof}
Use the family, $h$, $F$, $M$, and $B_h$ constructed above, and set
\begin{equation}\label{eq:polytope}
 K=\{x\in\R^d:F(x)=M\}
  =\bigcap_{\lambda\in\Lambda}
    \{x\in\R^d:r_\lambda\cdot x+c_\lambda(h)\ge M\}.
\end{equation}
The equality follows from $F\le M$. Thus $K$ is a closed convex
polyhedron containing $x^0$, and \eqref{eq:no-boundary-plateau} gives
$K\cap\partial B_h=\varnothing$.

In fact $K\subset\operatorname{int}B_h$. Otherwise a segment from
$x^0\in\operatorname{int}B_h$ to a point of $K$ outside that interior
would lie in $K$ by convexity and meet $\partial B_h$, a contradiction.
Therefore $K$ is bounded as well as closed, so it is a compact convex
polytope in $\R_{>0}^d$. This step uses no active-label estimate outside
the box.

Let $z\in K$. Since the polynomial vector field $f$ is smooth on
$\R^d$, it has a unique local solution $x$ from $z$ with a maximal
forward existence interval $[0,T_{\max})$. We initially regard this as
a solution in $\R^d$; confinement below will give positivity.
On any interval $[0,T]\subset[0,T_{\max})$ for which
$x([0,T])\subset B_h$, \eqref{eq:active-nondecrease} and
Lemma~\ref{lem:minimum-derivative} make $F(x(t))$ nondecreasing.
Since $F(z)=M$ and $F\le M$, we obtain
\begin{equation}\label{eq:trajectory-plateau}
 F(x(t))=M\qquad(0\le t\le T).
\end{equation}

The point $z$ lies in the interior of $B_h$. If the solution reached
the boundary at a first time $\tau<T_{\max}$, its whole segment
$x([0,\tau])$ would lie in $B_h$. Applying
\eqref{eq:trajectory-plateau} to that closed segment would give
$F(x(\tau))=M$, contrary to \eqref{eq:no-boundary-plateau}.
Continuity shows that any departure from the interior before
$T_{\max}$ would have such a first boundary time. Hence
$x(t)\in\operatorname{int}B_h$ for all $0\le t<T_{\max}$, and
\eqref{eq:trajectory-plateau} places the whole trajectory in $K$.

If $T_{\max}<\infty$, continuity of $f$ on the compact box gives
$L<\infty$ with $\norm{f(x)}\le L$ for $x\in B_h$. Consequently
\[
 \norm{x(t)-x(s)}\le L\abs{t-s}
 \qquad(0\le s,t<T_{\max}),
\]
so $x(t)$ has a limit $\bar x\in K\subset B_h$ as
$t\uparrow T_{\max}$. Local existence from $\bar x$ extends the
solution past $T_{\max}$: the original integral equation passes to
the endpoint, and both one-sided derivatives there equal $f(\bar x)$.
This contradicts maximality. Thus $T_{\max}=\infty$, proving the
claimed forward invariance for every $z\in K$.
\end{proof}

\begin{proof}[Proof of Theorem~\ref{thm:main}]
For the prescribed $x^0$, Proposition~\ref{prop:polytope} and its
construction give a forward-invariant polytope
$x^0\in K\subset B_h=[h,h^{-1}]^d$ with $0<h<1$. The solution
therefore exists globally and satisfies
\[
 h\le x_i(t)\le h^{-1}\qquad(t\ge0,\ 1\le i\le d).
\]
Taking $\varepsilon=h$ gives \eqref{eq:main-bounds}. Each coordinate
has positive infimum on the entire forward trajectory, and hence
positive forward limit inferior.
\end{proof}

\subsection{Equilibria and complex-balanced convergence}

The convexity of the invariant set also recovers the known existence of
positive equilibria in every positive stoichiometric class, proved by
Boros~\cite[preprint, Theorem~1]{Boros2019}.

\begin{corollary}\label{cor:positive-equilibrium}
For every finite weakly reversible mass-action network with positive
constant rates and every $x^0\in\R_{>0}^d$, there is an equilibrium
$x^*\in(x^0+S)\cap\R_{>0}^d$.
\end{corollary}

\begin{proof}
Let $K$ be the polytope in Proposition~\ref{prop:polytope}. The set
$Q=K\cap(x^0+S)$ is nonempty, compact, convex, and forward invariant.
Write $\phi_t$ for the flow. For each integer $n\ge1$, the continuous
map $\phi_{1/n}:Q\to Q$ has a fixed point $z_n$ by Brouwer's theorem.
Thus
\[
 0=n\bigl(\phi_{1/n}(z_n)-z_n\bigr)
   =n\int_0^{1/n}f(\phi_s(z_n))\,ds.
\]
Put $L=\max_{z\in Q}\norm{f(z)}$. Invariance and the integral equation
give $\norm{\phi_s(z_n)-z_n}\le L/n$ for $0\le s\le1/n$.
Uniform continuity of $f$ on $Q$ therefore yields
\[
 \norm{f(z_n)}
 \le\sup_{0\le s\le1/n}\norm{f(\phi_s(z_n))-f(z_n)}
 \longrightarrow0.
\]
A subsequence of $z_n$ converges to some $x^*\in Q$, and continuity gives
$f(x^*)=0$. Since $Q\subset\R_{>0}^d$, this is the required positive
equilibrium.
\end{proof}

\begin{remark}[Complex-balanced convergence]
\label{rem:complex-balanced-convergence}
Under the additional assumption of complex balance,
Theorem~\ref{thm:main} recovers the known global-attractor conclusion
(see also \cite{Craciun2026}): for fixed positive rate constants, every
positive trajectory converges to the unique positive equilibrium in its
positive stoichiometric class $P\cap\R_{>0}^d$.
Indeed, complex balance implies weak reversibility, and
Proposition~\ref{prop:polytope} traps the trajectory in the compact positive
forward-invariant set $K\cap(x^0+S)$. Let $x^*$ be the unique positive
equilibrium in this class. The Horn--Jackson entropy function
\[
 V(x)=\sum_{i=1}^d
 \left(x_i\log\frac{x_i}{x_i^*}-x_i+x_i^*\right)
\]
is nonincreasing and has zero derivative only at equilibria. This
dissipation property and uniqueness in each positive class are the
standard complex-balance conclusions
\cite[Proposition~5.3 and Corollaries~5.4--5.5]{FeinbergLectures1979}.
LaSalle's invariance principle therefore gives convergence.
In particular, the conclusion holds for every fixed positive rate choice
on every weakly reversible deficiency-zero network
\cite[Corollary~5.10]{FeinbergLectures1979}. Here the deficiency is
$|\mathcal C|-\ell-\dim S$, where $\ell$ is the number of linkage classes.
\end{remark}
